\section{Related Work}
\label{sec:related_work}

\paragraph{Semantic Correspondence.}
Early approaches to semantic correspondence focused on handcrafted features and geometric alignment techniques, which struggled to generalize
across large appearance and shape variations. More recent methods leverage deep learning to learn correspondence-aware representations,
significantly improving robustness to non-rigid deformations and intra-class variability. Benchmark datasets such as SPair-71k~\cite{min2019spair}
have played a crucial role in advancing this field by providing large-scale keypoint annotations across diverse object categories and viewpoints.

\paragraph{Emergent Correspondence from Pretrained Models.}
Recent studies have shown that semantic correspondence can emerge naturally from large pretrained visual models, even without explicit correspondence
supervision. Tang \etal~\cite{tang2023emergent} demonstrated that dense features extracted from image diffusion models encode strong correspondence
signals. Building on this idea, Zhang \etal~\cite{zhang2023tale} showed that features from Stable Diffusion complement self-supervised vision
transformers, enabling zero-shot semantic correspondence. These works highlight the surprising effectiveness of pretrained representations for dense
matching tasks and motivate the use of training-free baselines.

\paragraph{Self-Supervised Vision Transformers.}
Self-supervised learning has emerged as a powerful paradigm for learning general-purpose visual representations. DINO~\cite{caron2021emerging} showed
that Vision Transformers trained with self-distillation objectives capture high-level semantic structure and object part consistency. DINOv2~\cite{oquab2023dinov2}
further improved robustness and scalability by training on large curated datasets, leading to strong performance across a wide range of downstream
tasks. More recently, DINOv3~\cite{simeoni2025dinov3} has introduced architectural and training improvements that further enhance representation
quality. These models provide a strong foundation for semantic correspondence, as their dense features often align semantically meaningful object parts across images.

\paragraph{Foundation Models for Dense Vision Tasks.}
Beyond self-supervised encoders, foundation models trained for general vision tasks have shown strong transferability. Segment Anything~\cite{kirillov2023segment}
introduced a large-scale segmentation model capable of producing high-quality masks across diverse domains. Although not explicitly designed for
correspondence, its dense feature representations can be repurposed for matching tasks. Comparing such models with self-supervised transformers
provides insight into how different training objectives affect correspondence performance.

\paragraph{Prediction Strategies for Correspondence.}
Most correspondence pipelines rely on selecting the target location with the highest similarity score via a discrete argmax operation. However, this
approach is sensitive to noise and limited to pixel-level precision. Zhang \etal~\cite{zhang2024telling} proposed a geometry-aware framework that
replaces global argmax with a windowed soft-argmax strategy, enabling sub-pixel refinement and improved robustness. Inspired by this approach, we adopt
a similar prediction strategy to enhance correspondence accuracy in our experiments.

